\section{Discussion}
\label{sec:discussion}

The practical implication of our analysis is that not all approximation error
should be treated equally. For compression and sparse attention, an ideal
algorithm would preserve the direction-changing part of a residual update while
placing unavoidable error primarily in the parallel component. Error in this
component mostly rescales the current representation, whereas perpendicular
error changes its direction and is more closely associated with degradation in
our experiments. This suggests a geometric criterion for future compression
methods: in addition to total residual error, they should measure how much of
that error lies in the perpendicular component.

The same perspective also changes how we interpret training-time computation.
If value-space parallel manipulation has limited inference-time effect, then
repeatedly optimizing this self-value-parallel direction may not always be the
most useful allocation of model capacity or gradient signal. The retained
pretraining results in \S\ref{sec:pretraining} are consistent with this view:
self-value-parallel suppression changes optimization, and value-space parallel
removal gives the strongest average post-training improvement in the larger
retained summaries.
This does not imply that every residual-space parallel component is benign to
remove. Rather, it highlights the central lesson of the paper: the
importance of a parallel component depends on the space in which it is measured
and manipulated.
